\section{Related work}\label{sec:related}

\paragraph{Verifiable signature sharing.} Franklin and Reiter introduce the dealing of a finished signature among proxies who release it on a trigger, and the holder may be the signer~\cite{FR95}. Their ElGamal variant shares the signature scalar in base $r$ against a point that every proxy computes from the verification equation. Up to the sign convention this is the ECDSA sealed point. Their Schnorr variant pins the public image of the scalar by the challenge hash. Burmester gives an abstract form~\cite{Bur96}, and Catalano and Gennaro give a provably secure DSS variant that shares $s^{-1}$ in the base $g^m y^r$, the dual arrangement of ours~\cite{CG98}. These schemes send shares on private channels, share one signature per dealing, and leave the trigger to the judgment of the proxies. Camenisch and Damg{\aa}rd make signature sharing one publicly verifiable message with shares encrypted to each proxy~\cite{CD98}, and Fujisaki and Okamoto share an existing discrete-logarithm key with a proof against its public key~\cite{FO98}. In threshold proxy delegation, as the survey of Das et al.\ describes it, a Schnorr-type signature on a warrant is dealt with commitments whose constant term the verification equation forces~\cite{DSP06}. The result verifies under a derived proxy key, and is no ordinary signature of the owner.

\paragraph{Escrow, fair exchange and timed release.} Asokan, Shoup and Waidner verifiably escrow a Schnorr signature scalar to one trusted party that acts on dispute~\cite{ASW98}. Avoine and Vaudenay replace that party by arbiters who hold publicly verifiable encrypted shares and release them on a condition with an expiry~\cite{AV04}, and K{\"u}p{\c{c}}{\"u} and Lysyanskaya prove bounds for arbiters that never communicate~\cite{KL10}. In both lines the dealt item is generic and both exchanging parties stay active. Boneh, Gentry, Lynn and Shacham give verifiably encrypted BLS signatures whose check extends the verification equation~\cite{BGLS03}. The closest modern match to sealed mode is the verifiable timed signature of Thyagarajan et al.~\cite{VTS20}: a finished signature on a fixed transaction is Shamir-shared with its constant term forced by the verification equation, and the application includes conflicting pre-signed refund transactions on one output. There the release is by time and to anyone, no committee holds the shares or chooses among them, and every committed signature becomes public.

\paragraph{Dealt nonce material and linear signing.} Langford's threshold DSS with a trusted center is open-line dealing in its first form: for each future signature the center shares $k$ and $kx$, each shareholder answers with one linear partial, and a share must never be used twice, which the paper enforces by having every member acknowledge the message before anyone signs~\cite{Lan95}. It has no commitments, no public check of a partial, and its dealer is a key center. Gennaro, Jarecki, Krawczyk and Rabin observe that the public information of Feldman's scheme does not detect a wrong product share~\cite{GJKR96}. The check that does, in base $R$ with the sums of the two coefficient shares forced to $G$ and $X$, appears in the signing phase of Gennaro and Goldfeder~\cite{GG20} and of Canetti et al.~\cite{CGGMP21}, where the shares come from an interactive presigning protocol among the signers. Open mode joins the two: Langford's dealer with the base-$R$ check as a forced commitment on a public dealing. The Nillion report notes that a client could generate and distribute encrypted key shares centrally, with the nonce work left to the signers~\cite{Nillion23}. The service of Groth and Shoup states the one-use rule for its presignatures and re-randomizes each one after the request is fixed~\cite{GS22}. Open ECDSA does without re-randomization. It therefore rests on Theorem~3 of~\cite{GS21}, which needs ratio resistance and has the $UN$ term, where their service rests on the stronger Theorem~6 of that paper. For Schnorr, a dealer who Shamir-shares $x$ is the trusted-dealer setup of FROST~\cite{FROST,RFC9591}, and open Schnorr adds dealt nonce pairs for each slot.

\paragraph{Key import and owner-share protocols.} One-message publicly verifiable dealings can move a chosen key into a committee~\cite{cgVSS,Groth21}, and two-round ECDSA protocols accept a sharing made elsewhere~\cite{TroutPP,LLZD25}. None of these binds the dealing to an address or deals nonce material or signatures. The 2PC-MPC line keeps an owner share and lets the owner post one partial signature that the network completes when a condition holds~\cite{TwoPCMPC,FMMSS25,IkaDocs}. It is the nearest published match to the cryptographic one-spend limit of sealed mode, and it differs from sealed mode: the owner keeps a share, the owner's message consumes a presignature that the network made first, and one message authorizes one signature and no fee grid. Budgeted threshold signatures enforce a spending limit inside committee signing with an honest majority~\cite{CZC26}, which is a policy limit of the kind open mode has. The authorship proof of sealed ECDSA, a Schnorr proof of knowledge of the signature scalar in the base of the recovered nonce point, uses the rewriting of an ECDSA signature that zkAttest introduced and that a recent IETF draft specifies for device binding~\cite{zkAttestBlog,CLLZ26}. We add the binding of a dealing to an address.

\paragraph{Pre-signed transactions.} Vault designs replace a live key by pre-signed transactions and deleted keys, held and released by third parties on a chain event~\cite{SHMB20a,SHMB20b,SP21}, and they propose fee variants of one transaction as a hedge against the fee market~\cite{SHMB20a,SHMB20b,Todd23}. Covenant opcodes are the alternative that needs a consensus change~\cite{MES16,BIP345}; the motivation of BIP~345 lists the costs of pre-signed vaults, among them a fee policy that must be decided when the vault is created. Swambo et al.\ name verifiable secret sharing of the pre-signed transactions as future work and ask for a proof that a holder has a set of signatures without revealing them~\cite{SHMB20a}. The forced constant term of a sealed point is such a proof. Statechains and discreet log contracts hold conflicting pre-signed transactions on one outpoint and let the chain execute one~\cite{Mercury,Dryja}. These designs create a fresh key at deposit, and several rest on the honest deletion of a key or share, which \cref{thm:imposs} shows cannot bind a qualified holder.

\paragraph{Committee rules.} Quorum intersection in a committee of $3f+1$ is standard~\cite{Brick}. So is the anchor rule. Chainlink off-chain reporting collects $2f+1$ signed observations, sorts them and reports the median, which faulty oracles cannot move outside the range of the correct observations, and its pacemaker takes entry $f+1$ of the reported epochs~\cite{OCR21}. Deployed threshold systems gate signing on finalized external chain state~\cite{Chainflip,THORChain,TACo}. The Internet Computer minters and THORChain build transactions deterministically inside a replicated state machine, with a pinned account nonce or fixed inputs, a fee bound and replacement bumps~\cite{ckETH,ckBTC,THORChain}, and tBTC and Nomic chain their spends through one pinned output with a fee bound~\cite{tBTC,Nomic}. Chainflip, THORChain, tBTC and Nomic also move funds or keys across validator changes. All of them act for standing keys or scripts of the system itself, and none takes a dealing from an outside owner. What line dealing adds on this side is the application of these rules to a one-time intent bound to an outside owner's dealing: selection over a sealed grid, anchor retries that consume no nonce, reconstruction of earlier generations from board anchors alone, and a freeze that splits one request's dealt polynomials between two committees. Server-supported signatures enforce one signature per slot as a server policy with durable state~\cite{ATS96,DMT02}, and double-authentication-preventing signatures and online/offline signatures exhibit the two-points effect in other settings~\cite{PS17,ST01}.

\paragraph{Position.} We know of no work that combines an address-bound, publicly verifiable dealing by the owner of an existing address with committee selection or construction of one spend from finalized chain state. We do not claim priority for the dealing of committed presignatures: threshold ECDSA in the preprocessing model is a large literature, and such a dealer could appear in it.

\Cref{tab:forms} places the two modes beside the two published forms of \cref{sec:problem}.

\begin{table}[t]
\centering
\small
\begin{tabularx}{\linewidth}{@{}>{\raggedright\arraybackslash}p{2.9cm}>{\raggedright\arraybackslash}X>{\raggedright\arraybackslash}X>{\raggedright\arraybackslash}X>{\raggedright\arraybackslash}X@{}}
\toprule
 & Form A & Form B & Sealed & Open \\
\midrule
Existing address & separate dealing first & two owner transactions and a key registration & one message & one message \\
Owner reads committee state first & no & yes, a presignature & epoch keys only & epoch keys only \\
Raw key leaves the wallet & yes & yes, to split it & no & yes \\
Committee preprocessing & presignature protocol & network presigning & none & none \\
What a threshold coalition gets & the key, for all time & one stored partial signature & signatures on the owner's conflicting variants, never the key & the key \\
One-spend limit & policy & cryptographic, one message & cryptographic, one conflicting family & policy \\
Who builds the transaction & caller & owner or contract & owner builds variants, committee selects & committee \\
Failure before the spend & key stays with the committee & presignature consumed & coins never moved & coins never moved; nodes erase \\
\bottomrule
\end{tabularx}
\caption{The two modes of line dealing beside Form~A (the committee holds the key) and Form~B (the owner keeps a share).}
\label{tab:forms}
\end{table}
